% This fragment is intended as its own part. It requires only standard LaTeX
% plus amsmath/amssymb and hyperref/url for the bibliography.
\part{The findings in algebraic geometry and Grothendieck's frameworks}
\label{part:dictionary}

\section{Exact published and archival locators}
\label{sec:groth-locators}

The following locators identify the foundations carrying the
proofs across languages. They do not assert that Grothendieck's
texts state the modern tensor-complexity conclusions. PDF leaves
below are one-based counts in the specific facsimiles linked in
the bibliography; printed pagination is listed separately.

\begin{description}
\item[EGA properness.]
EGA~II \cite{EGA2}, Definition~5.4.1 and
Proposition~5.4.2(i)--(iv), printed p.~100, PDF leaf~97;
Theorem~5.5.3(i), printed p.~104, leaf~101.
These give universal closedness, the stability of proper
morphisms, and projective implies proper. Both leaves were
freshly rendered and visually checked in this run.

\item[FGA parameter spaces.]
Grothendieck's Bourbaki expos\'e~221 \cite{FGA},
Theorems~3.1--3.2, printed pp.~259--260, PDF leaves~12--13;
the Hilbert specialization after Proposition~3.8 and
Remarks~3.9, pp.~265--266, leaves~18--19.
These give the Quot/Hilbert parameter spaces and fixed-polynomial
projectivity used for limiting configurations. These four
facsimile leaves were freshly viewed.

\item[Group and module schemes.]
SGA~3 \cite{SGA3}, Expos\'e~I, \S\S4.5--4.6,
including Definition~4.6.1, printed p.~29, PDF leaf~42,
gives the general linear group and module functors. This leaf
was freshly viewed. In particular,
$\operatorname{GL}_n=\operatorname{Spec}
\mathbb Z[T_{ij},\det(T)^{-1}]$.

\item[Exterior powers, representation rings, and the gamma filtration.]
SGA~6 \cite{SGA6}, Expos\'e~V, Definition~2.1,
Examples~2.2(b), and Lemma~2.2.1, printed pp.~314--316,
PDF leaves~321--323; the representation-ring example before
\S2.3, p.~317, leaf~324; Proposition~3.2(ii), p.~325,
leaf~332; \S3.10, equation~(3.10.1), p.~331, leaf~338.
These supply the exterior-power operations on $K_0$, the
representation-ring example, and the gamma filtration.
All six listed leaves were freshly viewed.

\item[Scheme-theoretic images.]
EGA~I \cite{EGA1}, \S9.5, Propositions~9.5.1 and~9.5.4,
Corollary~9.5.2, Definition~9.5.3, printed pp.~176--177;
the predecessor records PDF leaves~175--176. This locator
is carried with its predecessor verification provenance,
not relabeled as a new facsimile check.

\item[Analytic tensor norms.]
The \emph{R\'esum\'e} \cite{Resume} is published-catalogue
item CSG17. The predecessor's inspected USP facsimile has
\S1 no.~1, Definition~1 and Theorem~1 at printed p.~8,
leaf~8; Theorem~2 and its duality corollary at p.~$8'$,
leaf~9. Its fundamental theorem, \S4 no.~2,
Theorems~1--2, is at pp.~59--60, leaves~60--61;
Corollary~1 after Theorem~3, inequality~(8), at p.~61,
leaf~62; the proof opening in \S4 no.~3 at p.~62,
leaf~63. These inherited locators are pertinent where
norms and duality are used, and are not a source for the
norm-free border-rank proof.

\item[\'Etale and weight terminology.]
SGA~4 \cite{SGA4}, Expos\'e~VII, \S1.2, supplies the
\'etale-site foundation. Deligne \cite[Definition~1.2.2,
p.~153]{WeilII} supplies purity terminology; the predecessor
records IAS PDF leaf~17. These source checks retain their
predecessor provenance. No SGA~5 theorem or Weil~II weight
theorem is used to deduce a tensor-rank bound here.
\end{description}

The primary CSG published-works catalogue \cite{CSGWorks}
was consulted in this run. CSG17 identifies the published
\emph{R\'esum\'e}, with its S\~ao Paulo volume and USP reprint;
CSG44 identifies the published Hilbert-schemes expos\'e and
its erratum. These identifiers are not Montpellier
unpublished-manuscript cotes. The Montpellier archive
\cite{MontpellierArchive} was consulted, but no manuscript
cote and scan carrying the modern findings below was
established. For every such finding the unpublished-manuscript
entry is therefore ``not established''; the exact published
foundations are the locators above. No cote or scan is invented.

\section{Dictionary for the residual-rank findings}

\paragraph{P-01: coordinate minima.}
The tensor factors are locally free sheaves over
$\operatorname{Spec}\mathbb Z$. The row/column rescalings
are characters of a split torus in the SGA~3 group schemes.
The residual image is a free submodule with its integral
coordinate basis. A $j$-plane is a point of the relative
Grassmannian. The echelon frames in P-01 give a rank-$j$
subbundle on $\mathbb A^1_k$: the pivot submatrix remains
the identity. Its tautological tensor specializes under
the product group action. The closed fiberwise Segre
secant cone contains the specialization whenever it
contains the nonzero-parameter fibers. This is exactly
the algebraic content of the proof; no metric is present.
The framework locators are SGA~3 Expos\'e~I \S\S4.5--4.6
and EGA~I \S9.5.

\paragraph{P-02: adjacent pairs.}
The maps identifying two distinct selected coordinate
sets are morphisms of free $\mathbb Z$-modules. Their
restrictions to each selected term are injective, and
the induced tensor maps send the residual generators
to the displayed adjacent generators with unit signs.
The Segre atom parametrization is functorial under these
maps. Hence the map sends every secant cone into the
corresponding target secant cone after each base change.
This is the sheaf/group-scheme restatement of the full
projection proof, with the same EGA~I and SGA~3 locators.
For $b>2$ this dictionary retains the exact adjacent-pair
reduction; it does not assign its unknown evaluated value.

\paragraph{P-03: stars.}
For $k\geq1$, the star certificate is an integral family
over $\mathbb A^1$: its displayed Laurent expression
cancels to a polynomial tensor. Over $\mathbb G_m$ it
is a sum of $k+1$ Segre sections. Its specialization
therefore lies in the fiberwise $(k+1)$-secant cone.
The lower bound is a unit minor of a locally free-sheaf
flattening. Both constructions commute with scalar
extension and survive every characteristic.
EGA~I \S9.5 supplies the image/closure language;
SGA~3 supplies the vector-bundle representation language.

\paragraph{P-04 and P-05: the exact pencil minima and Koszul certificate.}
The $p=1$ pencil map is the bundle morphism
\[
 A\otimes B^*\longrightarrow\Lambda^2 A\otimes C,
 \qquad a\otimes\beta\longmapsto
       \sum_h a_h\wedge a\otimes Q_h(\beta).
\]
Its matrix has integral entries. For the coordinate
terms its rank is the dimension of a union of coordinate
subbundles, certified by unit minors. A rank-one Segre
section has map rank one, so the relevant minors vanish
on the claimed lower secant cones. The upper family is
the sum of the explicit star families. This proves the
same exact values in the bundle language. SGA~6
Expos\'e~V, Examples~2.2(b) and Lemma~2.2.1, supplies
the exterior-power construction; the actual minors,
rather than only its $K_0$ class, supply the bound.

\paragraph{P-06: anchored shapes.}
The $h+r+s+1$ displayed vectors are sections of a free
module on $\mathbb A^1$. Their tensor squares are
sections of the complete Veronese cone in invariant
tensors. The four-term finite difference cancels its
constant and linear terms; after dividing by the
displayed power of the parameter, it extends regularly
to zero with the asserted residual tensor as value.
Grouping identical square sections gives the stated
number of Segre atoms on $\mathbb G_m$. The lower bound
is the common-kernel flattening and its unit minors.
This is an explicit secant degeneration in the EGA~I
image language. It requires no unproved smoothability
claim about an auxiliary finite scheme.

\paragraph{P-07: determinant compression.}
The row determinant image is the integral exterior-power
representation; the polynomial determinant image uses
the tensor product of the row and column exterior powers.
The lowering transvections are root subgroups of the
split Borel group scheme in the corresponding SGA~3
general linear group. Each two-element chain gives a
regular degeneration of Grassmannian frames, and the
integer descent terminates at a Borel-stable coordinate
subbundle. The proof uses the exterior-power cancellation
law, valid even in characteristic two, not semisimplicity.
SGA~6 Expos\'e~V, Examples~2.2(b), gives the relevant
exterior-power operation. The same unipotent stability
is not asserted for the permanent image in general.

\paragraph{P-08: finite Boolean differences.}
The vector $v(u,w,y,z)$ is a section of the row--column
tensor representation over the integral affine parameter
space. Its tensor square is a section of the complete
Veronese cone. Inclusion--exclusion on Boolean evaluation
sets extracts a squarefree coefficient after cancellation
of the indicated parameter powers. Grouping evaluations
gives a sum of eleven or twelve Segre sections over
$\mathbb G_m$ and a regular limiting section over
$\mathbb A^1$. This carries the full upper-bound proof
into EGA~I's image/closure framework and SGA~3's module
representation framework. It needs no division by an
integer and no smoothability assumption on an additional
scheme. Its conjunction with R-04 uses the latter's
properness and determinantal argument.

\paragraph{P-09, P-12, P-13, P-14 and P-16: seven-plane prolongation obstructions.}
The proper Grassmann incidence and exterior bundle maps are those of
R-04, with the exact EGA~II and SGA~6 locators recorded above.
A split torus (SGA~3, Expos\'e~I, \S\S4.5--4.6) acts on the
quotient of the tautological bundle. The finite certificates cover
all possible quotient weight supports, and then prove emptiness
of specified projective common-zero loci of linear maps.
P-14 and P-16 also enlarge the fixed subbundle before testing a
remaining variable line; the relevant annihilator is taken after
this enlargement. The map entries, their annihilators and their
minors carry the proof. Their source and target classes in $K_0$
alone do not. The modular and rational certificates have the
stated characteristic-zero scope. No unpublished-manuscript cote
for these modern computations has been established.

\paragraph{P-10 and P-11: integral curves and the row cubic seed.}
P-10's inclusion--exclusion expression is a sum of Segre sections over
$\mathbb G_m$ whose grouped sum extends over $\mathbb A^1$, with
integer ordered coefficients. P-11's six- and eight-term upper
curves are integral, without division by a nonunit integer.
Its double-Koszul map is an exterior-power morphism of locally
free sheaves. The determinant lower bound uses a proper product
of three Grassmannians, the closed multiplication-rank condition,
and a fixed-point argument for a triangular group. The stable
annihilator subbundles reduce to eight explicit triples whose
actual product-map minors are units times powers of two. This
explains exactly the exclusion of characteristic two in that
lower bound; there the determinant and permanent tensors coincide.
The foundations are EGA~I \S9.5, EGA~II \S5.4, SGA~3 Expos\'e~I
\S\S4.5--4.6, SGA~6 Expos\'e~V Definition~2.1, Examples~2.2(b)
and Lemma~2.2.1, and the FGA Grassmann/Hilbert/Quot framework at
the precise printed and facsimile locators above. These foundations
are not citations for the modern numerical values themselves.

\paragraph{P-15: exact parameter arrangements.}
After torus specialization, a quotient line is a point of a
projective weight space. Specifying the zero columns of an
induced morphism gives a locally closed stratum: a linear kernel
intersection with finitely many proper linear subspaces removed.
Over an infinite field this stratum is nonempty precisely when
none of those removed subspaces is the whole intersection.
Two-dimensional quotient subspaces in a three-dimensional weight
space are described by their annihilator lines in the dual
projective plane. The certificates exhaust these linear
arrangements, including compatibility after one quotient line
has been fixed. Thus the proof is an explicit emptiness proof
on projective Grassmann strata, using the same properness and
locally free morphisms as R-04. It does not replace a parameter
space by numerical samples or by a $K_0$ class.

\paragraph{P-17: Borel-stable extension flags.}
The necessary extension spaces form a closed locus in a proper
Grassmannian fiber. Its solvable-group fixed point is refined by
an invariant complete flag in the quotient of the tautological
bundle. Each flag step is a character line annihilated by the
adjacent root operators modulo the preceding subbundle. The
certificate computes these common kernels exactly and retains
any projective family as an unresolved parameter space. Thus
EGA~II properness, the SGA~3 integral module/group schemes and
FGA flag parameter spaces carry the proof. The eighteen named
exclusions follow from actual exterior-map minors; they are
not a consequence of a $K_0$ class or a claim about every
remaining ten-plane.

\paragraph{P-18: the polynomial cubic seed.}
The linear maps $\Phi$ and $\Psi$ between free tensor modules
transfer upper bounds through preimages of closed secant cones.
A $63$-minor of determinant $-16$ proves that the nine-labelled-square
map of affine schemes is dominant after inverting two: its
invertible linear term admits a formal inverse degree by degree.
The tangent curves then give the determinant upper bound twenty-seven;
characteristic two uses the integral permanent curve. The lower
map combines two exterior operations on locally free sheaves,
with explicit integer minors in its torus weight blocks.
These use the EGA~I closure, SGA~3 module-scheme and SGA~6 exterior
power locators already specified. The field conditions come from
the actual determinants and integral identities, not a categorical
dimension count or an assumed smoothability statement.

\paragraph{P-20: all row permanent quadratic minima.}
A graph containing a full star supplies an integral family of
$j+1$ labelled square tensors. First and mixed second differences
produce its quadratic edge inclusion tensor in the special fiber.
The lower bound is the proper substitution of R-01, using the
absence of a nonzero rank-one symmetric zero-diagonal matrix
over every field. The product of two such integral curves gives
the eighteen-term polynomial ten-plane upper bound. These are
explicit sections over $\mathbb G_m$ whose sum extends over
$\mathbb A^1$; the EGA image/closure and SGA~3 module-scheme
foundations apply after every base change. No additional
archival or motivic rank theorem is asserted.

\paragraph{P-19: a closed higher-prolongation condition.}
On the Grassmannian of candidate subbundles $E\subset A\otimes B$,
the condition
\[
 \dim\bigl((A\otimes E\otimes B)
       \cap(\Sym^2 A\otimes\Sym^2 B)\bigr)\geq r
\]
is a determinantal condition for a morphism out of a
constant-rank tautological bundle. The construction is made
in the fixed ambient bundle, so it remains closed even when
an intermediate first-prolongation kernel jumps in dimension.
Proper projection extends necessity from spans of independent
Segre points to their border limits. The successive kernel
calculations are then an exact way to evaluate this fixed
closed condition. This is EGA~II properness and the FGA
parameter-space framework applied to actual bundle morphisms;
the twenty-four final unit minors do not, by themselves,
change the field scope of the earlier rational arrangements.

\paragraph{R-01: multidimensional proper substitution.}
The image map is a regular map to a Grassmannian on
the open locus of constant flattening rank. A retained
$j$-plane is a point of its relative Grassmann bundle.
The low-border-rank condition on its tautological
tensor is closed in local frames and invariant under
changing frames. Proper projection of this closed
incidence gives the closed condition needed to pass
from actual decompositions to border decompositions.
The atom-span dimension argument is unchanged on
every algebraically closed fiber. These are precisely
EGA~II Definition~5.4.1, Proposition~5.4.2 and
Theorem~5.5.3(i), applied to the Grassmannian incidence;
FGA Theorems~3.1--3.2 provide the broader parameter-space
setting if limiting schemes are also retained.

\paragraph{R-02 and R-03: the evaluated higher-dimensional minima.}
The finite coordinate or Borel-ideal reductions parametrize
special points of the same relative Grassmannians.
Their ordinary and Koszul flattenings are explicit
morphisms of locally free sheaves. Nonzero minors give
the lower certificates; the anchored families give
the upper witnesses. This is the exact transfer of
the $j=3,4$ minima.
The SGA~3 actions, SGA~6 exterior powers, and EGA~II
closed-incidence locators above are the foundations.
Each certificate retains its stated characteristic
range. A determinant nonzero modulo one prime proves
nonzero over $\mathbb Q$, but does not by itself prove
nonzero in every characteristic; a unit minor does.

\paragraph{R-04: prolongation and the exceptional permanent subspace.}
The auxiliary limiting ten-plane is retained as a
Grassmannian point, together with the prescribed
subspace it contains. Its prolongation maps are
multiplication/intersection maps of tautological
bundles, so their rank conditions are determinantal
closed conditions. This permits passage to a limiting
span before applying the weight-block exclusions.
The actual exclusion matrices retain their entries
and minors; the recorded unit minors survive every
scalar extension. Properness is the EGA~II input.
FGA's Hilbert/Quot language can retain the limiting
configuration, but the proof does not identify its
actual span with the retained limiting plane when
an evaluation rank drops. The role of the embedding
and the actual multiplication maps is essential.

\paragraph{R-05: the fifth minimum and the bound for both polynomials.}
This deduction combines R-01's proper image condition,
the finite lower certificates including R-04, and
P-06's explicit upper family. In algebraic-geometric
terms it proves avoidance of the relevant lower
Grassmann secant locus, followed by avoidance of the
lower original Segre secant cone. It is asserted in
every characteristic: the permanent minors are units,
the determinant minors are units times powers of two,
and the coefficient tensors coincide in characteristic
two. It does not identify all
special fibers of an integral orbit or secant model
merely by reducing its characteristic-zero equations.

\paragraph{R-06: uniformity in the matrix size.}
Coordinate support submodules and their split inclusions
and projections are locally free-sheaf maps over
$\mathbb Z$. The degree-two support argument reduces
the relevant geometric points to a $4\times4$ coordinate
submodule or to five finite $5\times2$ graph orbits.
The ten further Koszul minors are integral determinantal
certificates. Their characteristic analysis is identical
to R-05, and the anchored upper curves are unchanged
under the split inclusions. The row assertion applies
the same proper Grassmann incidence inside the chosen
tautological residual subbundle and uses an integral
star curve. This is the proof in the EGA~II, SGA~3 and
SGA~6 frameworks with the exact locators already listed.

\paragraph{R-07: propagation along a flag.}
The pairs $V\subset W\subset U$ with dimensions $i,j$
form a relative flag variety, obtained by composing
Grassmann bundles. The proper-incidence argument for
R-01 applies on each fiber over $W$ and compares the
two tautological residual tensors. Taking minima gives
the monotonicity inequality; the complete anchored
$2\times3$ curve supplies the sixth minimum. Properness
is exactly EGA~II Proposition~5.4.2 and
Theorem~5.5.3(i). No new categorical invariant or
cohomological implication is needed for this deduction.

\paragraph{R-08, R-09 and R-10: complete covers and the bound forty-three.}
For the determinant the torus/root-group specialization of
P-07 reduces to a finite order-ideal cover. For the permanent,
row and column permutations and transposition supply a finite
symmetry cover after P-01; no permanent Borel invariance is
assumed. The independent certificates prove that every required
coordinate point belongs to a checked case. Exterior bundle
morphisms give the Koszul determinantal conditions, and proper
Grassmann incidences give the remaining prolongation conditions.
The anchored integral curves attain the minima. R-01 and its
flag version carry these minima back to the original tensor.
The foundational locators are the EGA~II, SGA~3, SGA~6 and FGA
locators already specified. The finite covers and their rank
conclusions are findings of this run, not theorems attributed
to Grothendieck's texts or to an unverified archive cote.

\paragraph{R-11: the uniform support-graph cover.}
Coordinate inclusions and split projections carry a residual
subbundle between ambient matrix sizes without changing border
rank. The support graph records only which free-module basis
sections occur. P-01 reduces to these coordinate subbundles;
the finite graph-and-subset cover exhausts them, and exterior
bundle morphisms with checked minors give the new rank bounds.
The underlying schemes and equivariant maps are those already
located in EGA~II, SGA~3 and SGA~6. The number of graph types is
a finite verified combinatorial calculation, not an invariant
recoverable from an unembedded motive or weight-zero stalk.

\paragraph{FL-01--FL-05: the all-field F-23 theorem.}
FL-01 defines integral tensor sections and makes explicit
the distinction between invariant tensors/divided powers
and the ordinary symmetric quotient. FL-02 is the
proper Grassmannian incidence above with $j=1$.
FL-03 identifies free residual image submodules and
split coordinate projections; unit coefficients
make the image bases survive every base change.
FL-04 is a unit-minor seed calculation. FL-05 is
the factorial obstruction in the ordinary symmetric
quotient. Together these are the complete dictionary
carrying the proofs, not an assumption that a
characteristic-zero normalization exists in every
fiber. Their precise published foundations are
EGA~II's properness statements, SGA~3's module and
group schemes, and SGA~6's locally free operations.

\paragraph{FL-06: the survey counterexample.}
The degree-two Veronese over $\mathbb C$ maps a vector
to its symmetric square. Its secant lower bound is
the rank of the first flattening, a determinantal
closed condition; the four-square expression is an
explicit point of the fourth secant parametrization.
The same proof is therefore an EGA~I image/closure
calculation on a locally free representation. No
additional archive or motivic input is used.

\section{The analytic findings and their precise boundary}

\paragraph{S-01: the exact DPS example.}
The underlying tensor and its three-term decomposition
are maps of free $\mathbb Z$-modules. Its residual
module is a rank-two direct summand of a free rank-six
module, hence a $\mathbb Z$-point of
$\operatorname{Gr}(2,6)$. The lower flattening bound
is an integral minor. These are the SGA~3 and EGA~I
parts. The projector and dual-witness matrices have
entries in $\mathbb Z[1/2]$, but their positive
semidefiniteness is a Hermitian statement on a chosen
complex fiber, certified by the displayed Gram
factorizations. The real inequality is not an
identity of algebraic $K_0$ classes.

\paragraph{S-02: the error estimate.}
The extension space uses the Schur functor
$\operatorname{Sym}^k B$ on a locally free sheaf.
Its rank is $\binom{k+\dim B-1}{\dim B-1}$, and its
$K_0$ generating operation is
$\sigma_t([B])=\lambda_{-t}([B])^{-1}$.
SGA~6 Expos\'e~V, Definition~2.1, Examples~2.2(b)
and Lemma~2.2.1, give this exterior/symmetric-operation
setting. The partial-transpose and positive-cone
constraints require a metric and a complex
archimedean fiber. The proved error bound is carried
by the analytic separability argument, not by a
gamma-filtration or a weight argument.

\paragraph{O-01: the orbit supremum.}
The product-$\operatorname{GL}$ group scheme acts on
the relative Grassmannian and on the Segre scheme.
The closed incidence
$\{U:\mathbb P(U)\cap\operatorname{Segre}\ne\varnothing\}$
is GCT\_GINEQ F-18's $\mathcal I_{q,1}$, obtained by
the EGA~II proper projection. The explicit
one-parameter family carries an accompanying line
to a Segre point; properness of the Grassmannian
and closed incidence retain that line in a limiting
subspace. The supremum-one conclusion is its
archimedean metric interpretation. For a product-free
space, the actual orbit misses this incidence while
its closure meets it. The diagonal action on an
alternating image is a different group action and
has the separate behavior stated in O-01.

\paragraph{O-02: the Pl\"ucker moment equation.}
The Pl\"ucker vector generates the determinant line
$\det(U)=\Lambda^q U$ of the tautological bundle,
inside $\Lambda^q(\bigotimes_i H_i)$. This exterior
power and group action are integral constructions
in SGA~3 and SGA~6. The Hermitian norm, partial
traces, and convexity along Hermitian one-parameter
subgroups are additional analytic structures.
The displayed scalar-marginal equation is the
critical-point condition for that norm, and its
proof establishes minimality. It is not an equality
of $K_0$ classes and is not the orbit-alpha supremum.

\section{What the categorical invariants do and do not retain}

For every Koszul map in this paper, $K_0$ records
the source and target classes, for example
$[B^*]\lambda^p[A]$ and $[C]\lambda^{p+1}[A]$.
For representations one uses the exact tensor
category of finite locally free representations;
SGA~6's representation-ring example is the relevant
locator. Proposition~3.2(ii) and equation~(3.10.1)
then describe the gamma operations and filtration.
Integral Schur functors can be used over
$\mathbb Z$; characteristic-zero Young projectors
requiring division by factorials or hook lengths
cannot be transferred without localization.

These constructions track modules and characters.
The rank of a particular morphism is supplied by
its minors, not by its source and target classes.
GCT\_GINEQ F-20 gives the precise counterexample:
the zero and identity maps have the same such
classes and different ranks. Hilbert and Quot
schemes retain additional geometric data, but
one must keep the actual embedding and limiting
span when those enter the certificate.

No new $\ell$-adic sheaf, purity theorem, or motive
is constructed by the proofs. GCT\_GINEQ F-21
already shows that the constant weight-zero stalk
and absolute point motive cannot distinguish a
zero tensor point from a nonzero pure tensor point.
This does not rule out embedding-sensitive sheaves
or relative motives as future invariants. It
identifies exactly why the present rank certificates
do not follow merely from attaching these names.
There is consequently no fabricated SGA~5 or
unpublished-manuscript theorem attached to such an
inference.

